\section{句法分析}
\label{sec:nlp-analysis-syntax}

词法分析给出了词和局部标签，却没有说明短语怎样嵌套，也没有说明哪个词支配哪个词。句法分析补上这两类信息。\term{成分结构}（constituency structure）用树表示相邻词怎样组成更大的短语，树的内部节点是短语类别；\term{依存结构}（dependency structure）用有向边表示中心词与依存词之间的关系，节点就是词本身。

图~\ref{fig:nlp-analysis-syntax}固定“小王／阅读／通知”的分词。成分树把“阅读通知”组成动词短语，再与“小王”组成句子；依存树以“阅读”为中心，分别连向主语“小王”和宾语“通知”。两种表示回答的问题不同：成分树说明哪些词连成一个单位，依存树说明词与词之间的支配关系。借助中心词规则可以把成分树转换成依存树，但反方向缺少短语类别和层次信息，两者一般不能无损互换。

\begin{figure*}[htbp]
  \centering
  \begin{tikzpicture}[
  x=1mm,y=1mm,
  every node/.append style={font=\figurefont\footnotesize,align=center},
  category/.style={draw=FigureModel,fill=FigureModel!10,
    rounded corners=1mm,inner sep=1.5mm},
  word/.style={draw=FigureInput,fill=FigureInput!10,
    minimum width=16mm,minimum height=8mm,inner sep=1mm},
  branch/.style={draw=NNLine,line width=.85pt},
  dep/.style={->,draw=NNWire,line width=1.1pt}
]
  \node[font=\figurefont\footnotesize\bfseries] at (40,59) {(a) 成分：短语包含关系};
  \node[category] (s) at (39,46) {S};
  \node[category] (np1) at (13,30) {NP};
  \node[category] (vp) at (53,30) {VP};
  \node[category] (v) at (41,15) {V};
  \node[category] (np2) at (68,15) {NP};
  \node[word] (a) at (13,0) {小王};
  \node[word] (b) at (41,0) {阅读};
  \node[word] (c) at (68,0) {通知};
  \foreach \parent/\child in {s/np1,s/vp,np1/a,vp/v,vp/np2,v/b,np2/c}
    \draw[branch] (\parent.south) -- (\child.north);
  \node at (13,-10) {$[0,2)$};
  \node at (41,-10) {$[2,4)$};
  \node at (68,-10) {$[4,6)$};

  \node[font=\figurefont\footnotesize\bfseries] at (127,59) {(b) 依存：词间有向关系};
  \node[category] (root) at (127,43) {0：人工根};
  \node[word] (d) at (102,0) {小王};
  \node[word] (e) at (127,0) {阅读};
  \node[word] (f) at (152,0) {通知};
  \draw[dep] (root.south) -- node[right,fill=white,inner sep=1mm] {root} (e.north);
  \draw[dep] (e.north west) to[out=145,in=45]
    node[above left,inner sep=1mm] {nsubj} (d.north);
  \draw[dep] (e.north east) to[out=35,in=135]
    node[above right,inner sep=1mm] {obj} (f.north);
  \node at (102,-10) {$1:[0,2)$};
  \node at (127,-10) {$2:[2,4)$};
  \node at (152,-10) {$3:[4,6)$};
  \draw[NNLine,dashed,line width=.5pt] (84,-14) -- (84,52);
\end{tikzpicture}

  \caption{同一句话的成分树与基本依存树。左图的NP、VP、S分别表示名词短语、动词短语和句子，方括号是字符跨度；右图的边由中心词指向依存词，0是人工根，nsubj和obj分别标记主语与宾语关系。两图固定同一分词与一种教学分析，不表示处理的先后。}
  \label{fig:nlp-analysis-syntax}
\end{figure*}

本节按输出结构分为两部分，每部分内部按第\ref{sec:nlp-analysis-paradigms}节的预测范式组织方法。成分解析依次覆盖区间动态规划（PCFG与CKY、Earley、神经跨度解析）、转移系统（移进--归约）和约束生成（序列化解析）；依存解析覆盖转移系统（arc-standard与arc-eager）和图上的精确解码（Eisner、Chu--Liu/Edmonds），再由双仿射网络提供弧分数，最后延伸到允许多入边的增强依存图。

\subsection{成分结构与解析}
\label{sec:nlp-analysis-constituency}

成分记为$(i,j,\ell)$，表示覆盖第$i$个词边界到第$j$个词边界之间的短语，类别为$\ell$。区间按词边界编号，不是字符偏移；输出时再经词跨度表还原原文。一棵树中的成分只能嵌套或互不相交，不能同时含有$[0,2)$和$[1,3)$这样交叉的区间。

“新／学生／代表”有两种分组：“新学生”的代表，即$[[\text{新 学生}]\ \text{代表}]$；或新任的“学生代表”，即$[\text{新}\ [\text{学生 代表}]]$。两种分组的词序完全相同，差别只在内部跨度。一个评分器可以给$[0,2)$与$[1,3)$都打高分，但合法树最多保留其中一个。这个例子集中体现了成分解析的四项挑战：一是\term{结构歧义}（structural ambiguity），同一词串对应多棵树，介词短语附着是英语中最常见的情形；二是上下文无关假设过强，规则概率只看当前类别，无法利用父节点和词汇信息；三是解析代价，在全部区间上搜索需要$O(n^3)$时间；四是输出必须是一棵合法树，局部最优的跨度拼不出合法结构。

方法谱系沿这几项挑战展开。PCFG与CKY提供可估计的概率树和精确搜索，父节点标注与词汇化缓解独立假设；Earley去掉乔姆斯基范式的要求，直接处理一般文法；神经跨度解析去掉显式文法，用整句上下文给跨度评分，仍以CKY保证合法；移进--归约把解析降到线性时间，代价是失去全局最优；序列化生成用统一的序列到序列接口输出括号树，合法性改由后处理或约束解码负责。

各方法用同一套指标比较。PARSEVAL评价把树转成带标签跨度的集合，计算预测集合与参考集合的精确率、召回率和F1，并事先固定是否计入根节点、词性节点和标点。以“新学生代表”为例，若参考为右分组、预测为左分组，计入根时参考集合为$\{\mathrm{NP}[0,3),\mathrm{NP}[1,3)\}$，预测集合为$\{\mathrm{NP}[0,3),\mathrm{NP}[0,2)\}$，只有根匹配，精确率、召回率和F1都是$1/2$；不计根时两者没有共同成分，F1为0。同一棵树在两种口径下的分数差别很大，比较不同系统时必须固定口径。

\subsubsection{PCFG与CKY图表解析}
\label{sec:nlp-analysis-pcfg-cky}

第\ref{sec:nlp-classical-pcfg-definition}小节已经定义了\term{概率上下文无关文法}（probabilistic context-free grammar, PCFG）：同一左部非终结符的展开概率之和为1，一棵树的概率是所用规则概率之积，见式~\eqref{eq:nlp-classical-pcfg}；第\ref{sec:nlp-classical-cky}小节给出了对所有树求和的inside算法，第\ref{sec:nlp-classical-inside-outside}小节给出了无树库时的inside--outside估计。本小节只处理解析：给定文法，求概率最大的树，并说明怎样从图表恢复原文中的成分。

有树库时，规则概率用相对频率估计：
\[
 p(A\to\beta)=\frac{N(A\to\beta)}{\sum_{\gamma}N(A\to\gamma)},
\]
$N(A\to\beta)$是规则在参考树中出现的次数；需要平滑时保持同一左部归一化。采用只有$A\to BC$与$A\to w$两类规则的乔姆斯基范式，本例固定根类别为NP，教学文法为
\[
\begin{array}{llll}
 \mathrm{NP}\to\mathrm A\ \mathrm{NP}&0.4&
 \mathrm{NP}\to\mathrm{NP}\ \mathrm N&0.3\\
 \mathrm{NP}\to\mathrm N\ \mathrm N&0.2&
 \mathrm{NP}\to\text{学生}&0.1\\
 \mathrm A\to\text{新}&1&
 \mathrm N\to\text{学生}&0.6\\
 &&\mathrm N\to\text{代表}&0.4
\end{array}
\]
NP规则概率之和为1，N规则概率之和也为1。这是只用于手算的微型文法；没有对应规则的组合不是低概率树，而是根本不在候选集合中。

\term{CKY图表解析}（CKY chart parsing）把每个“区间加根类别”作为可复用的状态。令$V[A,i,j]$为类别$A$覆盖$[i,j)$时最优子树的概率。长度为1的区间直接查词汇规则，$V[A,i,i+1]=p(A\to w_{i+1})$，没有规则时为0。长度至少为2时，根若使用$A\to BC$，左右子树必在某个边界$k$处分开，于是
\begin{equation}
 V[A,i,j]=
 \max_{\substack{A\to BC\\i<k<j}}
 p(A\to BC)\,V[B,i,k]\,V[C,k,j].
 \label{eq:nlp-analysis-pcfg-cky}
\end{equation}
这就是inside递推式~\eqref{eq:nlp-classical-inside}把求和换成取最大的形式。按区间长度从小到大填表，每个组合所需的两个子区间都已算好；同时保存取得最大值的规则与分割点$k$，供回溯使用。最优子结构成立的理由是：若最优树的某棵子树不是同一类别、同一区间下的最优子树，把它换成最优子树会提高整棵树的概率，与最优性矛盾。

表~\ref{tab:nlp-analysis-cky}列出所有非零条目。长度为2的“学生代表”有两条推导，必须先比较它们，而不是保留第一次找到的规则。

\begin{table}[htbp]
  \centering
  \begin{tabularx}{\linewidth}{@{}l X l@{}}
    \toprule
    区间&非零类别及计算&最优回溯\\
    \midrule
    $[0,1)$&$\mathrm A:1$&新\\
    $[1,2)$&$\mathrm N:0.6,\ \mathrm{NP}:0.1$&学生\\
    $[2,3)$&$\mathrm N:0.4$&代表\\
    $[0,2)$&$\mathrm{NP}:0.4\cdot1\cdot0.1=0.04$&$\mathrm A,\mathrm{NP},k=1$\\
    $[1,3)$&$\mathrm{NP}:\max(0.2\cdot0.6\cdot0.4,\,
      0.3\cdot0.1\cdot0.4)=0.048$&$\mathrm N,\mathrm N,k=2$\\
    $[0,3)$&$\mathrm{NP}:\max(0.4\cdot1\cdot0.048,\,
      0.3\cdot0.04\cdot0.4)=0.0192$&$\mathrm A,\mathrm{NP},k=1$\\
    \bottomrule
  \end{tabularx}
  \caption{“新／学生／代表”的完整非零CKY图表。未列出的条目为0。根的另一种顶层分组概率为0.0048；右侧区间$[1,3)$内部另有概率0.012的推导$\mathrm{NP}\to\mathrm{NP}\ \mathrm N$，但最优值为0.048。}
  \label{tab:nlp-analysis-cky}
\end{table}

从根$V[\mathrm{NP},0,3]$读出$(\mathrm A,\mathrm{NP},k=1)$，再读出右子树的$(\mathrm N,\mathrm N,k=2)$，得到
\[
 [\mathrm{NP}\ [\mathrm A\ \text{新}]
       [\mathrm{NP}\ [\mathrm N\ \text{学生}]
                    [\mathrm N\ \text{代表}]]],
\]
概率为0.0192，左分组为$0.3\cdot0.04\cdot0.4=0.0048$。原文“新学生代表”中三个词的字符跨度依次为$[0,1)$、$[1,3)$、$[3,5)$，所以右侧NP还原为“学生代表”$[1,5)$。这个结论只反映本例文法的偏好；若前文是“刚入学”，读者可能偏向左分组，而PCFG读不到这条背景。

实际文法并不都是二叉的。规则$A\to BCD$可以引入专用辅助类别$A_\star$，改写成$A\to B\,A_\star$与$A_\star\to CD$，原概率放在第一条，第二条设为1，回溯后删去辅助节点；辅助类别必须记住原规则身份，否则会引入原文法不存在的组合。单分支链$A\to B\to C$可以折叠成复合类别，恢复时再展开；空产生式需要事先消除。若根值为0，返回解析失败。长句用对数概率把乘积变成和，避免下溢。设二叉规则集合为$R_2$、非终结符集合为$N$，直接实现的时间为$O(n^3|R_2|)$，图表与回溯指针的空间为$O(n^2|N|)$。

微型文法的另一个问题是NP在不同位置共用一个分布。第\ref{sec:nlp-classical-annotation}小节介绍的\term{父节点标注}（parent annotation）把类别改写成“类别加父类别”，让展开概率依赖结构位置。以下计数是教学设定：设树库中100个NP节点的展开与上面的概率一致，其中60个是句根NP，40个位于另一个NP之下。按父类别拆开后，
\[
\begin{array}{lcccc}
 &\mathrm A\ \mathrm{NP}&\mathrm{NP}\ \mathrm N&\mathrm N\ \mathrm N&\text{学生}\\
 \mathrm{NP}^{\wedge}\mathrm{ROOT}\ (60)&18&30&12&0\\
 \mathrm{NP}^{\wedge}\mathrm{NP}\ (40)&22&0&8&10
\end{array}
\]
各列相加仍是$40,30,20,10$。重新估计后，右分组的概率为$p(\mathrm{NP}^{\wedge}\mathrm{ROOT}\to\mathrm A\ \mathrm{NP}^{\wedge}\mathrm{NP})\cdot p(\mathrm{NP}^{\wedge}\mathrm{NP}\to\mathrm N\ \mathrm N)\cdot0.6\cdot0.4=0.3\cdot0.2\cdot0.24=0.0144$，左分组为$0.5\cdot(0.55\cdot0.25)\cdot0.4=0.0275$，偏好翻转。原文法把两种位置的统计平均在一起，掩盖了“句根NP常以中心名词结尾、内嵌NP常是单个名词”这类差异；标注后类别数增加，每条规则的训练样本变少，需要平滑。Klein与Manning沿这一思路系统地拆分类别，不用词汇化，在英语宾州树库第23节40词以内的句子上得到86.36的带标签F1，接近同期的词汇化解析器，说明“按什么状态共享展开概率”本身就决定了解析能力\scite{v9a-klein2003}。

\subsubsection{Earley文法解析}
\label{sec:nlp-analysis-earley}

CKY要求文法是乔姆斯基范式。文法作者未必愿意把每条规则改成二叉形式，二叉化还会改变规则的含义和概率的放置位置。\term{Earley解析}（Earley parsing）不改写文法，用带点的规则记录“一条规则已经读到了哪里”\scite{v9a-earley1970}。

\term{带点项目}（dotted item）$[A\to\alpha\cdot\beta,i]$存放在图表位置$j$上，表示这条规则从边界$i$开始，已经匹配了$w_{i+1:j}$对应的$\alpha$，还在等待$\beta$。点是进度标记，不是输入词元。图表位置$j$的项目集合记作$C_j$。为了让句根也用同一机制处理，增加一条起始规则$S'\to\mathrm{NP}$。

\begin{definition}[Earley的三条推理规则]
  \label{def:nlp-analysis-earley-rules}
  初始时$[S'\to\cdot\,\mathrm{NP},0]\in C_0$。设输入为$w_1\cdots w_n$，按以下规则向图表加入项目，同一位置上相同的项目只保留一份：
  \[
  \begin{array}{lll}
   \text{预测}&
   \dfrac{[A\to\alpha\cdot B\beta,\,i]\in C_j\qquad B\to\gamma\in G}
         {[B\to\cdot\,\gamma,\,j]\in C_j}\\[3.2ex]
   \text{扫描}&
   \dfrac{[A\to\alpha\cdot a\beta,\,i]\in C_j\qquad w_{j+1}=a}
         {[A\to\alpha a\cdot\beta,\,i]\in C_{j+1}}\\[3.2ex]
   \text{完成}&
   \dfrac{[B\to\gamma\,\cdot,\,k]\in C_j\qquad[A\to\alpha\cdot B\beta,\,i]\in C_k}
         {[A\to\alpha B\cdot\beta,\,i]\in C_j}
  \end{array}
  \]
  横线上方是前提，下方是结论。若$[S'\to\mathrm{NP}\,\cdot,0]\in C_n$，则输入被文法接受。
\end{definition}

预测说明“等待$B$的项目需要从当前位置开始识别一个$B$”；扫描让点越过与下一个输入词相同的终结符，并把项目移到下一个位置；完成把一个已经识别完的$B$交还给所有在其起点等待$B$的项目。每个位置反复应用预测与完成直到不再产生新项目，再对下一个词做扫描。

沿用上一小节的文法，位置0预测出NP的四条规则$[\mathrm{NP}\to\cdot\,\mathrm A\ \mathrm{NP},0]$、$[\mathrm{NP}\to\cdot\,\mathrm{NP}\ \mathrm N,0]$、$[\mathrm{NP}\to\cdot\,\mathrm N\ \mathrm N,0]$、$[\mathrm{NP}\to\cdot\,\text{学生},0]$，以及A、N的词汇规则。规则$\mathrm{NP}\to\mathrm{NP}\ \mathrm N$是左递归的：预测它又会预测NP，但这些项目在$C_0$中已经存在，去重使预测停止。这正是Earley能直接处理左递归、而朴素的自顶向下递归下降会无限循环的原因。

扫描“新”后，$C_1$得到$[\mathrm A\to\text{新}\,\cdot,0]$；完成操作回到位置0，推进出$[\mathrm{NP}\to\mathrm A\cdot\mathrm{NP},0]$，后者又预测从1开始的NP。扫描“学生”后，$C_2$同时有$[\mathrm N\to\text{学生}\,\cdot,1]$与$[\mathrm{NP}\to\text{学生}\,\cdot,1]$：前者推进$[\mathrm{NP}\to\mathrm N\cdot\mathrm N,1]$；后者完成$[\mathrm{NP}\to\mathrm A\ \mathrm{NP}\,\cdot,0]$，即“新学生”，并推进$[\mathrm{NP}\to\mathrm{NP}\cdot\mathrm N,1]$。由于“新学生”是一个完整的NP，它又推进$[\mathrm{NP}\to\mathrm{NP}\cdot\mathrm N,0]$。扫描“代表”后，$C_3$完成三条NP：$[\mathrm{NP}\to\mathrm N\ \mathrm N\,\cdot,1]$与$[\mathrm{NP}\to\mathrm{NP}\ \mathrm N\,\cdot,1]$都覆盖“学生代表”，$[\mathrm{NP}\to\mathrm{NP}\ \mathrm N\,\cdot,0]$覆盖整句的左分组；前两者再完成$[\mathrm{NP}\to\mathrm A\ \mathrm{NP}\,\cdot,0]$，即整句的右分组。最终$[S'\to\mathrm{NP}\,\cdot,0]\in C_3$，输入被接受。

接受项目只说明至少存在一棵树。要输出树，每次扫描和完成都要记录结论由哪些前提项目得到；同一项目可以有多组来源，这些来源指针构成\term{共享解析森林}（shared packed forest）。本例中$[\mathrm{NP},0,3]$有两组来源，$[\mathrm{NP},1,3]$也有两组。给项目附上Viterbi内部概率：预测时设为规则概率，完成时乘上被交还子项目的概率，同一项目的多组来源取最大，就得到与CKY相同的结果：右分组0.0192胜过左分组0.0048。Stolcke把这种做法推广为同时计算前缀概率的概率Earley解析\scite{nlp-analysis-stolcke1995}。

可空规则$B\to\varepsilon$需要单独处理。$B$在位置$j$预测后立即完成，若此时等待$B$的项目尚未全部加入$C_j$，完成就会漏掉它们。Aycock与Horspool的修正是：预测一个可空类别时，同时把等待它的项目的点直接越过它\scite{nlp-analysis-aycock2002}。单分支循环（如$A\to B$、$B\to A$）会使森林含环，Viterbi计算需要在环上迭代或预先消除。Earley在一般文法下的最坏时间为$O(n^3)$，在无歧义文法上为$O(n^2)$，对许多确定性文法接近线性。它与CKY的主要区别在于不必改写文法，以及只在可由左侧推出的位置上预测；共同代价是歧义越多，森林越大。

\subsubsection{神经跨度成分解析}
\label{sec:nlp-analysis-neural-cky}

父节点标注和词汇化都在手工设计“规则概率应当依赖什么”。神经跨度解析换一条路：去掉显式文法，用编码整句得到的上下文表示直接给每个跨度和类别评分，再用CKY保证输出是合法树。它属于第\ref{sec:nlp-analysis-paradigm-span}小节的跨度范式，评分对象是带标签的跨度$(i,j,\ell)$。

跨度表示的一般形式是式~\eqref{eq:nlp-analysis-span-representation}的端点拼接、内部注意力汇聚与宽度嵌入。成分解析中常用一种更紧凑的形式。用双向LSTM在句首、句末加哨兵符号后编码整句，在边界位置$t=0,\ldots,n$上取前向状态$\bm f_t\in\mathbb R^d$与后向状态$\bm b_t\in\mathbb R^d$：$\bm f_t$读完了边界$t$左侧的词，$\bm b_t$读完了右侧的词。跨度$(i,j)$的表示定义为
\begin{equation}
 \bm r_{i,j}=[\bm f_j-\bm f_i;\ \bm b_i-\bm b_j]\in\mathbb R^{2d}.
 \label{eq:nlp-analysis-span-difference}
\end{equation}
$\bm f_j$与$\bm f_i$的差近似刻画前向网络读入$w_{i+1:j}$这几个词带来的变化，$\bm b_i-\bm b_j$是后向网络读同一段时的变化，因此不需要额外的宽度嵌入或内部汇聚。图~\ref{fig:nlp-analysis-span-difference}画出了“学生代表”$(1,3)$的四个状态。这种差分是Wang与Chang在依存解析中提出的LSTM-Minus特征的双向形式；Cross与Huang把它用于转移式成分解析，Stern等用于图表解析\cite{nlp-analysis-cross2016,nlp-analysis-stern2017}。Kitaev与Klein把编码器换成自注意力，把每个位置的输出向量拆成两半分别充当$\bm f_t$与$\bm b_t$，差分形式不变\scite{nlp-intro-kitaev2018}。

\begin{figure*}[htbp]
  \centering
  \begin{tikzpicture}[
  x=1mm,y=1mm,figure picture,
  word/.style={figure node,minimum height=7mm,minimum width=16mm,
    draw=FigureInput,fill=FigureInputFill},
  post/.style={circle,draw=FigureStroke,fill=white,line width=.75pt,
    minimum size=5.5mm,inner sep=0pt,font=\figurefont\scriptsize},
  state/.style={figure node,minimum height=6mm,minimum width=11mm,
    inner xsep=1pt,font=\figurefont\scriptsize},
  hot/.style={state,draw=FigureOutput,fill=FigureOutputFill,line width=1pt},
  fwd/.style={figure flow},
  bwd/.style={figure flow,dash pattern=on 2pt off 1.5pt},
  note/.style={font=\figurefont\scriptsize,text=FigureMuted}
]
  % 词与边界位置
  \foreach \k/\w in {1/新,2/学生,3/代表}{
    \node[word] (w\k) at (\k*30-15,0) {\w};
  }
  \foreach \k in {0,1,2,3}{
    \node[post] (p\k) at (\k*30,-11) {\k};
  }
  \node[note,anchor=east] at (-6,-11) {边界位置};

  % 前向与后向状态
  \foreach \k in {0,1,2,3}{
    \node[state] (f\k) at (\k*30,13) {$\bm f_{\k}$};
    \node[state] (b\k) at (\k*30,25) {$\bm b_{\k}$};
  }
  \node[hot] at (f1) {$\bm f_{1}$};
  \node[hot] at (f3) {$\bm f_{3}$};
  \node[hot] at (b1) {$\bm b_{1}$};
  \node[hot] at (b3) {$\bm b_{3}$};
  \foreach \a/\b in {0/1,1/2,2/3}{
    \draw[fwd] (f\a) -- (f\b);
    \draw[bwd] (b\b) -- (b\a);
  }
  \node[note,anchor=east] at (-6,13) {前向（读完左侧）};
  \node[note,anchor=east] at (-6,25) {后向（读完右侧）};

  % 待评分跨度
  \draw[figure subgroup,draw=FigureOutput] (24.5,-17) rectangle (95.5,6);
  \node[note,anchor=north] at (60,-17.5) {跨度 $(1,3)$：学生代表};

  % 差分表示：四个高亮状态参与差分
  \node[figure node,draw=FigureOutput,fill=white,minimum width=96mm,
    minimum height=8mm] (r) at (45,40)
    {$\bm r_{1,3}=[\bm f_{3}-\bm f_{1};\ \bm b_{1}-\bm b_{3}]\in\mathbb R^{2d}$};
  \foreach \k in {1,3}{\draw[figure link] (b\k.north) -- (b\k.north |- r.south);}
  \node[note,anchor=south] at (45,44.5) {再经 MLP 得到各标签分数 $s(1,3,\ell)$};
\end{tikzpicture}

  \caption{边界差分的跨度表示。下排圆圈是词之间的边界位置0至3；中排前向状态$\bm f_t$已读完边界$t$左侧的词，上排后向状态$\bm b_t$已读完右侧的词。跨度$(1,3)$只用高亮的四个状态：前向差$\bm f_3-\bm f_1$与后向差$\bm b_1-\bm b_3$都对应读入“学生代表”带来的变化。}
  \label{fig:nlp-analysis-span-difference}
\end{figure*}

标签分数由一个单隐层前馈网络给出，这里采用Kitaev与Klein加入层归一化的写法：
\begin{equation}
 \bm s_{i,j}=\bm W_2
  \operatorname{ReLU}\!\bigl(
   \operatorname{LayerNorm}(\bm W_1\bm r_{i,j}+\bm b_1)
  \bigr)+\bm b_2\in\mathbb R^{L},
 \label{eq:nlp-analysis-span-label}
\end{equation}
其中$\bm W_1\in\mathbb R^{d_h\times2d}$，$\bm W_2\in\mathbb R^{L\times d_h}$，$L$是标签数，$\bm s_{i,j}$的第$\ell$个分量记为$s(i,j,\ell)$。一句话共有$n(n+1)/2$个跨度，全部分数可组织成形状为$(n+1)\times(n+1)\times L$、只用上三角的张量，一次批量计算。

树的分数是所含跨度分数之和，$S(T)=\sum_{(i,j,\ell)\in T}s(i,j,\ell)$。单分支链折叠成一个复合标签，例如S与VP的链记作S+VP；二叉化引入的辅助节点使用空标签$\varnothing$，固定$s(i,j,\varnothing)=0$。这样同一棵原始树的不同二叉化不会仅因辅助节点而得到不同总分；空标签的跨度在恢复时删去。若标签选择只依赖跨度本身，令$q(i,j)=\max_\ell s(i,j,\ell)$，最优树由式~\eqref{eq:nlp-analysis-span-cky}的跨度CKY求出。在本节的边界编号下，$C(i,j)$表示区间$[i,j)$上最优子树的总分，递推为
\begin{align}
 C(i,i+1)&=q(i,i+1),\nonumber\\
 C(i,j)&=q(i,j)+\max_{i<k<j}[C(i,k)+C(k,j)].
 \label{eq:nlp-analysis-neural-cky}
\end{align}
状态里不再有类别，复杂度从$O(n^3|R_2|)$降到$O(n^3)$加上$O(n^2L)$的标签取最大。若评分含父子类别耦合，就须像式~\eqref{eq:nlp-analysis-pcfg-cky}那样把类别留在状态中。

对同一“新／学生／代表”，设单词叶分数为0，两个长度为2的跨度的最佳NP分数分别为5和6，根NP为1，其他非空类别都低于这些值（教学数值）。独立选择高分跨度会同时取$[0,2)$与$[1,3)$，总分$5+6+1=12$，但二者交叉。图表计算$C(0,2)=5$、$C(1,3)=6$，根比较$1+0+6=7$与$1+5+0=6$，选择右分组。

训练使用式~\eqref{eq:nlp-analysis-structured-objectives}中的结构化间隔损失，具体为
\begin{equation}
 \mathcal L=\max_T[S(T)+\Delta(T,T^\star)]-S(T^\star),
 \label{eq:nlp-analysis-parse-hinge}
\end{equation}
$T^\star$是参考树，$\Delta(T,T^\star)$是$T$中不属于$T^\star$的带标签跨度个数。候选包含$T^\star$本身，所以损失非负。$\Delta$按跨度分解，代价增强解码只需把每个跨度分数改为$s(i,j,\ell)+\mathbb 1[(i,j,\ell)\notin T^\star]$再跑一遍CKY，就能找到“最违反间隔”的树；梯度提高参考跨度的分数、降低这棵竞争树中错误跨度的分数。损失为0意味着参考树比每棵其他树至少高出它们的错误跨度数。

与PCFG相比，神经跨度解析的收益来自整句上下文：跨度分数可以利用远处的词和字符形态，不再受规则局部独立的限制。Kitaev与Klein沿用同一图表解码、把编码器换成自注意力，不用外部预训练表示时单模型在宾州树库WSJ测试集上的F1为93.55\scite{nlp-intro-kitaev2018}。代价是要有树库监督，且精确解码依赖分数按跨度分解；不同跨度之间的相互作用无法进入这个CKY。

\subsubsection{移进--归约成分解析}
\label{sec:nlp-analysis-shift-reduce}

图表方法同时保存$O(n^2)$个区间的分数，解码需要$O(n^3)$时间。长句或流式输入时，可以改用第\ref{sec:nlp-analysis-paradigm-transition}小节的转移系统，按动作逐步构树。移进--归约成分解析和下一小节的arc-standard依存解析属于同一范式：状态由栈、缓冲区和已建结构组成，每个动作有前提与效果，参考结构经oracle转成动作序列，模型在合法动作上做softmax，按式~\eqref{eq:nlp-analysis-transition-loss}以交叉熵训练。两者的区别只在动作集合与栈上存放的对象：这里栈上是带类别的子树，依存解析的栈上是词。

Sagae与Lavie、Zhu等的成分系统使用表~\ref{tab:nlp-analysis-shift-reduce-actions}中的动作\scite{nlp-analysis-sagae2005,nlp-analysis-zhu2013}。归约动作中的L或R指明新节点的中心词取自左孩子还是右孩子，这是为了让特征能读取子树的中心词；X是新节点的类别。参考树需要先二叉化，并记录单分支的位置。

\begin{table*}[htbp]
  \centering
  \begin{tabularx}{\linewidth}{@{}l X X@{}}
    \toprule
    动作&前提&效果\\
    \midrule
    SHIFT&缓冲区非空&把缓冲区首词连同其词性作为叶节点压栈\\
    REDUCE-L-X / REDUCE-R-X&栈中至少两个子树，且合成后满足二叉化约束&弹出栈顶两个子树，合成类别为X的节点后压栈；中心词取自左（L）或右（R）孩子\\
    UNARY-X&栈非空，且栈顶刚才没有执行过单分支&弹出栈顶子树，套一层类别为X的节点后压栈\\
    FINISH&缓冲区为空，栈中只剩一个可作根的子树&终止，输出栈中的树\\
    \bottomrule
  \end{tabularx}
  \caption{移进--归约成分解析的动作。前提保证每一步都可执行，且不会出现无限单分支；它们只保证结构合法，不保证选中的分组正确。}
  \label{tab:nlp-analysis-shift-reduce-actions}
\end{table*}

对右分组，“新／学生／代表”的参考轨迹为五个构树动作加终止：
\[
\begin{array}{llll}
 \text{步}&\text{动作}&\text{栈}&\text{缓冲区}\\
 1&\mathrm{SHIFT}&[\mathrm A(\text{新})]&[\text{学生},\text{代表}]\\
 2&\mathrm{SHIFT}&[\mathrm A,\ \mathrm N(\text{学生})]&[\text{代表}]\\
 3&\mathrm{SHIFT}&[\mathrm A,\ \mathrm N,\ \mathrm N(\text{代表})]&[\,]\\
 4&\mathrm{REDUCE\mbox{-}R\mbox{-}NP}&[\mathrm A,\ \mathrm{NP}(\text{学生代表})]&[\,]\\
 5&\mathrm{REDUCE\mbox{-}R\mbox{-}NP}&[\mathrm{NP}(\text{新学生代表})]&[\,]
\end{array}
\]
随后FINISH。左分组的轨迹则在第2步之后立即归约：SHIFT、SHIFT、REDUCE-R-NP（新学生）、SHIFT、REDUCE-R-NP。两条轨迹在第3步分叉：读完“新／学生”之后，选择继续移进还是立即归约，决定了最终分组。这个选择是不可撤销的；一旦把“新学生”归约成NP，后续任何动作都不能再把它拆开。

动作分类器读取栈顶若干子树的类别与中心词、缓冲区前几个词等信息。设分叉处“继续移进”与“立即归约”的局部概率为$.45$与$.55$，两条后续路径的条件概率分别为$.9$与$.4$，完整轨迹概率为$.405$与$.22$（教学数值）。贪心解码在分叉处选择归约，得到左分组；束宽为2的束搜索保留两条，最终按完整概率选回右分组。束搜索仍是近似，有限束宽不能保证找到全局最优。单分支使不同树的动作序列长度不同，比较束中候选时会偏向较短的轨迹；Zhu等为此加入空动作IDLE，让所有候选在相同步数后比较。动态oracle与全局归一化等通用的训练改进见第\ref{sec:nlp-analysis-paradigm-transition}小节。

移进--归约的时间与句长成线性，每步只读取有限的状态特征，适合需要快速解析的场景。代价是早期错误无法撤销，且没有CKY那样的全局最优保证。Cross与Huang的转移式成分解析把动作分类与上一小节的跨度差分结合：动作分类器读取栈顶跨度的差分表示，并给出可证明最优的动态oracle\scite{nlp-analysis-cross2016}。

\subsubsection{序列化生成式解析}
\label{sec:nlp-analysis-seq2seq-parsing}

前三种方法都为树设计了专门的搜索。另一条路线完全不设计搜索：把树写成括号序列，交给通用的序列到序列模型生成，属于第\ref{sec:nlp-analysis-paradigm-generation}小节的约束生成范式。Vinyals等按深度优先顺序线性化树，词本身不出现在输出中，叶节点上的词性标签统一替换为占位符XX，例如右分组写成
\[
 \texttt{(NP XX (NP XX XX )NP )NP},
\]
右括号带上类别，便于模型记住当前要关闭哪一层\scite{nlp-analysis-vinyals2015grammar}。编码器读入原句，带注意力的LSTM解码器逐个生成括号与XX，训练目标是参考序列每一步的交叉熵。推断用束搜索；解码完成后检查括号是否配对、XX的个数是否等于词数，不合法时补全或删除括号，再把第$k$个XX对应回第$k$个词。原文中，只用WSJ训练时无注意力的模型F1不到70，加入注意力后才可用；再用大量自动解析的句子扩充训练集，达到了与专门解析器相当的水平。

这条路线的收益是接口统一：同一个序列到序列模型可以解析、翻译或生成摘要，不需要为树设计特征和动态规划。代价是合法性不再由构造保证。括号配对是上下文无关语言的性质，可以用第\ref{sec:nlp-inference-grammar}小节的文法约束解码在每步屏蔽不能完成合法括号串的词元，从源头排除不配对的输出；叶节点个数与原句一致则需要在约束中计数。

大语言模型可以在提示中给出括号格式和少量示例，直接输出成分树。这时还会出现一种序列化解析特有的错误：模型在括号中改写、漏掉或增加原句的词。交付前应逐一核对叶序列与原句的词序列，并用PARSEVAL在固定口径上评价；没有树库微调时，模型输出的类别体系可能与目标树库不一致，需要映射。

表~\ref{tab:nlp-analysis-constituency-methods}在连续成分树这一共同输出下比较五条路线。

\begin{table*}[htbp]
  \centering
  \begin{tabularx}{\linewidth}{@{}l XXXX@{}}
    \toprule
    方法&解决的问题&核心机制&推断与复杂度&新代价\\
    \midrule
    PCFG--CKY&概率化歧义选择&规则概率之积；区间动态规划&精确，$O(n^3|R_2|)$&独立假设强；需二叉化\\
    Earley&一般文法不必改写&带点项目；预测、扫描、完成&精确，最坏$O(n^3)$&可空规则与森林处理复杂\\
    神经跨度--CKY&利用整句上下文&边界差分表示；跨度标签分数&精确，$O(n^3)$&需树库；分数须按跨度分解\\
    移进--归约&降低解析时间&栈上子树与归约动作&近似（贪心或束），$O(n)$&早期错误不可撤销\\
    序列化生成&统一序列接口&括号线性化；seq2seq生成&近似（束），$O(n)$步&合法性靠约束或后处理\\
    \bottomrule
  \end{tabularx}
  \caption{成分解析方法对照。复杂度中的$n$为句长；输出均须去掉辅助节点、展开复合类别后，再按同一PARSEVAL口径评价。}
  \label{tab:nlp-analysis-constituency-methods}
\end{table*}

\subsection{依存结构与解析}
\label{sec:nlp-analysis-dependency}

基本依存分析以词$1,\ldots,n$和人工根0为节点，边$h\to d$表示中心词$h$支配依存词$d$，边上另有关系标签。本节要求：每个真实词恰有一个中心词，人工根没有入边，所有真实词可从根到达，并且人工根恰好连接一个真实词。满足这些条件的边集合是以0为根的有向生成树。\term{投射性}（projectivity）进一步要求每个词及其全部后代在句中占据连续区间；把人工根放在最左侧时，等价于画在句子上方的依存边互不交叉。依存标注资源以通用依存（Universal Dependencies, UD）为代表，它为多种语言提供共享的关系标签体系\scite{nivre2020universal}。

评价按词计分。\term{无标签依存准确率}（unlabeled attachment score, UAS）是中心词预测正确的词所占比例，\term{带标签依存准确率}（labeled attachment score, LAS）还要求关系标签正确。本节手算计入指向根词的根弧、示例不含标点；实际评测需声明是否排除标点，并在不同系统间保持分母一致。

依存解析的挑战有四项：输出必须满足单头、单根、无环的树约束；自由语序和长距离依赖会产生非投射结构；关系标签需要与弧一起预测；长距离弧的错误率明显高于短弧。方法谱系为：转移式解析线性时间但贪心；图式解析把树的分数分解到弧上，在一阶分解下全局精确，其中Eisner算法限于投射树，Chu--Liu/Edmonds允许非投射；双仿射网络用神经弧评分代替手工特征，接图解码；增强依存图取消单头限制，表达共享论元与省略。每种方法都要说明它在什么结构空间内保证合法。

\subsubsection{arc-standard与arc-eager转移解析}
\label{sec:nlp-analysis-arc-transition}

两种转移系统都维护栈$\sigma$、缓冲区$\beta$和弧集合$A$，从$([0],[1,\ldots,n],\varnothing)$开始，到缓冲区为空、栈中只剩根0时结束\scite{nlp-analysis-nivre2003,nlp-analysis-nivre2008}。它们的差别在于连弧的时机，表~\ref{tab:nlp-analysis-arc-systems}列出两套动作。记栈顶为$s_0$、次顶为$s_1$，缓冲区首词为$b_0$。

\begin{table*}[htbp]
  \centering
  \begin{tabularx}{\linewidth}{@{}l l X X@{}}
    \toprule
    系统&动作&前提&效果\\
    \midrule
    arc-standard&SHIFT&缓冲区非空&$b_0$压栈\\
    &LEFT-ARC-$r$&栈中至少两个元素，$s_1\neq0$&加弧$s_0\xrightarrow{r}s_1$，移除$s_1$\\
    &RIGHT-ARC-$r$&栈中至少两个元素；$s_1=0$时缓冲区须为空&加弧$s_1\xrightarrow{r}s_0$，移除$s_0$\\
    \midrule
    arc-eager&SHIFT&缓冲区非空&$b_0$压栈\\
    &LEFT-ARC-$r$&$s_0\neq0$，且$s_0$尚无中心词&加弧$b_0\xrightarrow{r}s_0$，弹出$s_0$\\
    &RIGHT-ARC-$r$&缓冲区非空；$s_0=0$时根尚无孩子&加弧$s_0\xrightarrow{r}b_0$，$b_0$压栈\\
    &REDUCE&$s_0$已有中心词&弹出$s_0$\\
    \bottomrule
  \end{tabularx}
  \caption{两种依存转移系统的动作。arc-standard只在栈顶两个元素之间连弧，被连上的依存词立即移除，因此必须等它的孩子全部连完；arc-eager在栈顶与缓冲区首词之间连弧，右依存词一出现就连上，之后用REDUCE清理。两套系统都只能产生投射树。}
  \label{tab:nlp-analysis-arc-systems}
\end{table*}

对“小王$_1$／阅读$_2$／通知$_3$”，arc-standard的完整轨迹为：SHIFT，SHIFT，LEFT-ARC-nsubj建立$2\to1$，SHIFT，RIGHT-ARC-obj建立$2\to3$，RIGHT-ARC-root建立$0\to2$。栈依次为
\[
 [0]\to[0,1]\to[0,1,2]\to[0,2]
 \to[0,2,3]\to[0,2]\to[0].
\]
根弧在最后才建立，此时“阅读”的子树已经完整。arc-eager的轨迹为：SHIFT，LEFT-ARC-nsubj建立$2\to1$，RIGHT-ARC-root建立$0\to2$，RIGHT-ARC-obj建立$2\to3$，REDUCE，REDUCE。它在读入“阅读”时就连上了根，比arc-standard早；这种“尽早连弧”让arc-eager的部分结构更早可用，也让它成为最早得到动态oracle的系统之一，Goldberg与Nivre证明了它的弧可以逐条判断能否仍被到达\scite{nlp-analysis-goldberg2012}。

两套系统各有$n$次移进和$n$次连弧，arc-eager另有若干次REDUCE，动作数都与句长成线性。没有额外动作时它们只能产生投射树：被移除或弹出的词不再回到栈中，交叉弧所需的连接顺序无法实现。Nivre在arc-standard中加入SWAP动作，把次栈顶的词移回缓冲区，从而改变词进入栈的顺序；按重排后的顺序连出的投射弧，在原词序中可以交叉，解析时间在期望意义上仍是线性的\scite{nlp-analysis-nivre2009}。

\paragraph{Chen与Manning的前馈动作分类器}
早期转移解析器用人工组合的稀疏特征给动作打分，特征模板多、组合稀疏，特征提取占去大部分解析时间。Chen与Manning把离散特征换成稠密嵌入，用一个小型前馈网络预测arc-standard的动作\scite{v9a-chen2014}。输入取三组元素：栈顶三个词与缓冲区前三个词；栈顶两个词的最左、最右两个孩子；以及这两个词的最左孩子的最左孩子、最右孩子的最右孩子。词特征共18个，对应的词性特征18个，除栈与缓冲区6个词之外的12个元素再取弧标签，共48个特征。每个特征查表得到50维嵌入，按组拼接成$\bm x^{w}\in\mathbb R^{18\cdot50}$、$\bm x^{t}\in\mathbb R^{18\cdot50}$、$\bm x^{l}\in\mathbb R^{12\cdot50}$，隐藏层与输出层为
\begin{equation}
 \bm h=\bigl(\bm W_1^{w}\bm x^{w}+\bm W_1^{t}\bm x^{t}+\bm W_1^{l}\bm x^{l}+\bm b_1\bigr)^{3},\qquad
 p(a\mid c)=\operatorname{softmax}(\bm W_2\bm h)_a,
 \label{eq:nlp-analysis-chen-manning}
\end{equation}
立方逐元素计算，$\bm h\in\mathbb R^{200}$，$\bm W_2$的行数是动作数$2N_r+1$（$N_r$为关系标签数）。立方激活展开后含有三个输入分量的乘积项，相当于自动学习特征两两、三三组合，替代了人工组合模板。训练以oracle动作的交叉熵加$L_2$正则，嵌入用预训练词向量初始化。实现上，作者对高频词预先计算嵌入与$\bm W_1$对应列的乘积并缓存，解析时只需查表相加。在宾州树库（CoNLL依存转换）的开发集上，它以92.2的UAS每秒解析约一千句，准确率与速度同时高于作者实现的稀疏特征转移解析器。

\paragraph{栈LSTM}
Chen与Manning的特征仍是固定窗口：栈中更深处的词和已归约子树的内部结构看不到。Dyer等用\term{栈LSTM}（stack LSTM）编码整个栈\scite{nlp-analysis-dyer2015}。栈LSTM是带栈顶指针的LSTM：压栈时以当前栈顶的隐状态为前一状态计算新状态，弹栈时只把指针退回上一个状态，因此栈顶的输出始终概括了栈中全部元素。解析器用三个这样的LSTM分别编码栈$\sigma$、缓冲区$\beta$和动作历史，状态表示为
\[
 \bm p_t=\max\{\bm 0,\ \bm W[\bm s_t;\bm b_t;\bm a_t]+\bm d\},
\]
再经softmax给出动作分布。归约时，中心词$\bm h$、依存词$\bm d$与关系嵌入$\bm e_r$组合为$\bm c=\tanh(\bm U[\bm h;\bm d;\bm e_r]+\bm e)$，作为新子树压回栈中。与固定窗口相比，栈LSTM能读到完整的解析历史，代价是每步都要更新递归状态。

转移解析的共同限制是贪心：每步只看局部动作分数，一旦错误弹出某个词，涉及它的后续弧就不可达。动态oracle、束搜索与全局归一化的改进见第\ref{sec:nlp-analysis-paradigm-transition}小节。

\subsubsection{Eisner投射树解码}
\label{sec:nlp-analysis-eisner}

图式解析先给每条候选弧$h\to d$打分$s(h,d)$，再在合法树中找总分最大的一棵，$\max_{T}\sum_{(h,d)\in T}s(h,d)$。这种一阶分解是第\ref{sec:nlp-analysis-paradigm-graph}小节图范式的典型例子。逐词选最高入边可能成环；即使打断环，结果也可能不投射。Eisner算法在连续区间上递推，直接在投射树空间内求精确最优\scite{nlp-analysis-eisner1996}。

Eisner算法的四类状态与递推已在第\ref{sec:nlp-analysis-paradigm-graph}小节式~\eqref{eq:nlp-analysis-eisner}给出，本小节沿用其记号：闭区间$[i,j]$表示从词$i$到词$j$，与成分解析的边界区间相区别；完整项$C$的根在区间一端，它在区间内一侧的子孙都已收齐；不完整项$I$刚加入端点弧，被支配的端点还要收齐朝向区间内的孩子。实现时把非初值条目设为$-\infty$，按区间长度从短到长、同一长度先算不完整项后算完整项，并记录每次取最大值的分割点和状态类型，供回溯使用。本小节补上图范式中没有展开的两件事：单根约束和一个完整的手算例子。

单根约束在图表完成之后单独处理。一棵投射树被根词$r$分成左右两半：左半是以$r$为右端根的完整区间$C^{\leftarrow}[1,r]$，右半是以$r$为左端根的完整区间$C^{\to}[r,n]$。因此最优值为
\[
 \max_{1\leq r\leq n}
 \bigl[s(0,r)+C^{\leftarrow}[1,r]+C^{\to}[r,n]\bigr],
\]
并保存最优的根词$r$。若把人工根当作位置0直接放进图表递推，它可能连接多个真实词，得到的是一片森林挂在根下，不满足单根要求。

给“小王／阅读／通知”构造全部有限的弧分数（教学数值）：
\[
 s(0,2)=5,\quad s(2,1)=4,\quad s(2,3)=3,\quad
 s(3,1)=5,\quad s(1,3)=6,
\]
其余弧禁止，包括自环和指向人工根的弧。逐词取最高入边得到$3\to1$、$0\to2$、$1\to3$，总分16，但1与3之间成环。长度为2的图表项为$I^{\leftarrow}[1,2]=C^{\leftarrow}[1,2]=4$与$I^{\to}[2,3]=C^{\to}[2,3]=3$，相反方向均为$-\infty$。弧$1\to3$虽有6分，但$I^{\to}[1,3]$的两个分割分别需要$C^{\leftarrow}[2,3]$与$C^{\to}[1,2]$，即被禁止的$3\to2$或$1\to2$，这一不完整项不可达。唯一可行的根是2，总分$5+4+3=12$；回溯得到$2\to1$、$2\to3$与$0\to2$，正是图~\ref{fig:nlp-analysis-syntax}的无标签树。

Eisner算法的时间为$O(n^3)$，图表与回溯指针为$O(n^2)$。它本身不学习参数：弧分数可以来自结构感知器或大间隔训练\cite{mcdonald-etal-2005-non}，训练时用当前分数解码出竞争树，再比较参考树与竞争树的弧特征；也可以来自下文双仿射网络的局部交叉熵。精确性只针对一阶可加的弧分数和投射树空间；加入兄弟、祖父等高阶部件后，状态要扩大，复杂度随之上升。

\subsubsection{Chu--Liu/Edmonds非投射树解码}
\label{sec:nlp-analysis-mst}

自由语序和远距离依赖会产生交叉弧。取消投射性之后，解码问题变成在带权有向图上求以0为根的\term{最大生成树形图}（maximum spanning arborescence）：每个真实词恰有一条入边，无环且可从根到达。McDonald等把非投射依存解析写成这一问题\cite{v9a-mcdonald2005}，求解用\term{Chu--Liu/Edmonds算法}（Chu--Liu/Edmonds algorithm）\scite{nlp-analysis-chu1965,nlp-analysis-edmonds1967}。

算法的思路是：先让每个真实词贪心地选最大入边；若没有环，这已经是上界，直接返回。若有环，最优树一定恰好从环外进入环一次，于是把环收缩成一个点，进入这个点的弧要反映“替换掉环中哪条边”的代价。设环$Q$中节点$v$的已选前驱为$a(v)$，环内总分$w(Q)=\sum_{v\in Q}s(a(v),v)$。外部弧$u\to v$进入环时替换掉原入边$a(v)\to v$，收缩后的分数为
\begin{equation}
 s'(u,Q)=\max_{v\in Q}
   [s(u,v)-s(a(v),v)+w(Q)].
 \label{eq:nlp-analysis-mst-contract}
\end{equation}
出环弧取$s'(Q,u)=\max_{v\in Q}s(v,u)$。这就是第\ref{sec:nlp-analysis-paradigm-graph}小节给出的收缩规则，只是这里把常数$w(Q)$也加回，使收缩图上的树分直接等于原图上展开后的树分。两种最大值都要保存原始端点，否则展开时无法恢复词之间的弧。算法~\ref{alg:nlp-analysis-cle}给出递归形式。

\begin{algorithm}[Chu--Liu/Edmonds最大生成树形图]
  \label{alg:nlp-analysis-cle}
  \begin{algorithmic}[1]
    \Require 节点集合$V=\{0,1,\ldots,n\}$，弧分数$s(u,v)$（禁止的弧为$-\infty$），根0
    \Ensure 以0为根、总分最大的生成树形图$T$
    \Procedure{CLE}{$V,s$}
      \For{每个$v\in V\setminus\{0\}$}
        \State $a(v)\gets\argmax_{u\neq v}s(u,v)$\Comment{贪心选最大入边}
      \EndFor
      \State $M\gets\{(a(v),v):v\neq0\}$
      \If{$M$中无环}
        \State \Return $M$
      \EndIf
      \State 任取$M$中的一个环$Q$，计算$w(Q)=\sum_{v\in Q}s(a(v),v)$
      \State 把$Q$收缩为新节点$q$，得到$V'=(V\setminus Q)\cup\{q\}$
      \For{每个$u\notin Q$}
        \State $s'(u,q)\gets\max_{v\in Q}[s(u,v)-s(a(v),v)]+w(Q)$，记录取到最大值的$v^\ast(u)$
        \State $s'(q,u)\gets\max_{v\in Q}s(v,u)$，记录取到最大值的源点$v^{\circ}(u)$
      \EndFor
      \State 环外节点之间的弧分数不变；$T'\gets$\Call{CLE}{$V',s'$}
      \State 设$T'$中进入$q$的弧为$(u_0,q)$：加入$u_0\to v^\ast(u_0)$，并加入$Q$中除$a(v^\ast(u_0))\to v^\ast(u_0)$以外的全部弧
      \State 把$T'$中每条$(q,u)$换回$v^{\circ}(u)\to u$，其余弧原样保留
      \State \Return 展开后的弧集合$T$
    \EndProcedure
  \end{algorithmic}
\end{algorithm}

沿用上一组分数，环$Q=\{1,3\}$的总分为$5+6=11$。从2进入1的收缩分数为$4-5+11=10$，进入3的为$3-6+11=8$，从0进入的弧都被禁止。收缩图中$q$的最大入边是$2\to q$（10分），2的最大入边是$0\to2$（5分），无环，递归结束。展开时用$2\to1$替换环边$3\to1$，保留$1\to3$，得到总分$5+4+6=15$的树：$0\to2$、$2\to1$、$1\to3$。这棵树无环，但弧$0\to2$与$1\to3$交叉，因而不投射；Eisner排除了它，最高只能得到12分。两者的分差完全来自允许的树空间不同，不说明15分的那棵树在语言上更合理。

如第\ref{sec:nlp-analysis-paradigm-graph}小节所述，朴素实现每次收缩花$O(n^2)$、最多收缩$n$次，总时间$O(n^3)$，Tarjan的实现在稠密图上为$O(n^2)$。上述算法允许人工根连接多个真实词。若规范只允许一个根词，最直接的做法是枚举根词$r$，只保留$0\to r$这一条根出边，分别求解后取最高分，代价是多出$n$倍的计算；Zmigrod等给出了在$O(n^2)$内直接施加单根约束的做法\scite{nlp-analysis-zmigrod2020}。

非投射解码并非总是更好。McDonald等在捷克语PDT上比较两种解码，生成树解码的中心词准确率为84.4，投射解码为83.3；而在依存树全部投射的英语宾州树库转换数据上，投射解码反而略高\scite{v9a-mcdonald2005}。捷克语中非投射弧不足2\%，收益却已可见；在没有交叉弧的数据上，更大的树空间就没有可恢复的对象。选哪一种应看目标树库中非投射结构的比例。

\subsubsection{双仿射依存解析}
\label{sec:nlp-analysis-biaffine}

上面两个解码器都假定弧分数已经给定。Chen与Manning的分类器给的是动作分数，不是弧分数；早期图式解析器用人工特征给弧打分。Kiperwasser与Goldberg先用BiLSTM的输出向量代替人工特征，再用前馈网络给每对词打分\scite{nlp-analysis-kiperwasser2016}。Dozat与Manning的\term{双仿射依存解析}（biaffine dependency parsing）进一步改变评分函数：先把每个词分别投影成“作为依存词”和“作为中心词”的两种表示，再用双仿射形式一次算出全部$n\times n$的弧分数\scite{nlp-analysis-dozat2017}。图~\ref{fig:nlp-analysis-biaffine}画出了整体结构。

\begin{figure*}[htbp]
  \centering
  \begin{tikzpicture}[
  x=1mm,y=1mm,figure picture,
  input/.style={figure node,minimum height=7mm,minimum width=22mm,
    draw=FigureInput,fill=FigureInputFill},
  layer/.style={figure node,minimum height=8mm},
  mlp/.style={figure node,minimum height=8mm,minimum width=25mm},
  hot/.style={figure node,minimum height=9mm,draw=FigureOutput,
    fill=FigureOutputFill,line width=1pt},
  flow/.style={figure flow},
  shape note/.style={font=\figurefont\scriptsize,text=FigureMuted}
]
  % 输入与编码
  \node[input] (x) at (60,0) {词与词性嵌入 $\bm x_0,\ldots,\bm x_n$};
  \node[layer,minimum width=132mm] (enc) at (60,14) {多层 BiLSTM：$\bm R\in\mathbb R^{(n+1)\times 2h}$};
  \draw[flow] (x) -- (enc);

  % 四个投影
  \node[mlp] (ad) at (8,31) {$\mathrm{MLP}^{\mathrm{arc\mbox{-}dep}}$};
  \node[mlp] (ah) at (42,31) {$\mathrm{MLP}^{\mathrm{arc\mbox{-}head}}$};
  \node[mlp] (rd) at (78,31) {$\mathrm{MLP}^{\mathrm{rel\mbox{-}dep}}$};
  \node[mlp] (rh) at (112,31) {$\mathrm{MLP}^{\mathrm{rel\mbox{-}head}}$};
  \foreach \m in {ad,ah,rd,rh}{\draw[flow] (enc.north -| \m.south) -- (\m.south);}

  % 两个双仿射头
  \node[hot,minimum width=58mm] (sa) at (25,52)
    {$\bm S^{\mathrm{arc}}=\bm H^{(h)}\bm U\bm H^{(d)\top}+(\bm H^{(h)}\bm u)\bm 1^{\top}$};
  \node[hot,minimum width=58mm] (sr) at (95,52)
    {关系双仿射：$\bm S^{\mathrm{rel}}\in\mathbb R^{(n+1)\times(n+1)\times R}$};
  \draw[flow] (ad.north) -- (ad.north |- sa.south);
  \draw[flow] (ah.north) -- (ah.north |- sa.south);
  \draw[flow] (rd.north) -- (rd.north |- sr.south);
  \draw[flow] (rh.north) -- (rh.north |- sr.south);

  % 解码
  \node[layer,minimum width=46mm] (tree) at (25,72) {Eisner 或 CLE 选树};
  \node[layer,minimum width=46mm] (lab) at (95,72) {逐弧取最高关系};
  \draw[flow] (sa) -- (tree);
  \draw[flow] (sr) -- (lab);
  \draw[flow] (tree) -- node[above,shape note] {弧集合} (lab);
\end{tikzpicture}

  \caption{双仿射依存解析器的结构。BiLSTM输出经四个独立的MLP投影为弧依存端、弧中心端（维度$d_a$）与关系依存端、关系中心端（维度$d_r$）表示；弧头用双仿射得到$(n+1)\times(n+1)$的弧分数矩阵并交给树解码器，关系头只在选中的弧上读取$R$个关系分数。行下标为中心词，列下标为依存词，第0行对应人工根。}
  \label{fig:nlp-analysis-biaffine}
\end{figure*}

设句子连同人工根共$n+1$个位置，多层BiLSTM输出$\bm R\in\mathbb R^{(n+1)\times2h}$，第$i$行$\bm r_i$是位置$i$的上下文表示。两个MLP把它们投影到较低的维度$d_a$：
\[
 \bm h^{(d)}_i=\mathrm{MLP}^{\mathrm{arc\mbox{-}dep}}(\bm r_i),\qquad
 \bm h^{(h)}_i=\mathrm{MLP}^{\mathrm{arc\mbox{-}head}}(\bm r_i),
\]
按行堆叠成$\bm H^{(d)},\bm H^{(h)}\in\mathbb R^{(n+1)\times d_a}$。弧分数矩阵为
\begin{equation}
 \bm S^{\mathrm{arc}}=\bm H^{(h)}\bm U\bm H^{(d)\top}+(\bm H^{(h)}\bm u)\bm 1^{\top}\in\mathbb R^{(n+1)\times(n+1)},
 \label{eq:nlp-analysis-biaffine}
\end{equation}
其中$\bm U\in\mathbb R^{d_a\times d_a}$，$\bm u\in\mathbb R^{d_a}$，$\bm 1$是全1列向量。元素$S^{\mathrm{arc}}_{i,j}=\bm h^{(h)\top}_i\bm U\bm h^{(d)}_j+\bm h^{(h)\top}_i\bm u$是“$i$作$j$的中心词”的分数。双线性项衡量两端是否相容；第二项只依赖中心词，表示某个词本身有多大可能带依存词，例如动词比限定词更常作中心词。依存端自身的偏置对同一列的所有候选中心词相同，不影响选择，所以省去。

关系标签用另一对MLP投影到维度$d_r$，得到$\bm g^{(d)}_j$与$\bm g^{(h)}_i$，对每个关系$r$计算
\[
 S^{\mathrm{rel}}_{i,j,r}=\bm g^{(h)\top}_i\bm U^{(r)}\bm g^{(d)}_j+\bm w_r^{\top}[\bm g^{(h)}_i;\bm g^{(d)}_j]+b_r,
\]
$\bm U^{(r)}\in\mathbb R^{d_r\times d_r}$。把$R$个$\bm U^{(r)}$叠成三阶张量，全部分数构成$(n+1)\times(n+1)\times R$的张量。原文的BiLSTM为三层、每方向400维，弧投影$d_a=500$，关系投影$d_r=100$。

为什么要先投影再配对，而不直接用$\bm r_i$？BiLSTM的输出要同时服务于弧、关系和其他信息，维度高且混杂；专用MLP把对当前判断有用的部分抽出，并把双线性项的参数从$(2h)^2$降到$d_a^2$。弧与关系使用不同的投影，是因为“谁支配谁”与“是什么关系”依赖的信息不同。

训练时，对每个依存词$j$在第$j$列上做softmax，以参考中心词$i^\star$训练$-\log p(i^\star\mid j)$；关系头只在参考弧上训练正确关系的交叉熵。推断时先把$\bm S^{\mathrm{arc}}$交给Eisner或Chu--Liu/Edmonds选出合法树，再在选中的每条弧上取最高的关系分数。逐列取最大虽然也能给出每个词的中心词，但可能成环或多根；树解码器保证输出合法。

取$d_a=2$、$\bm U=\bm I$、$\bm u=\bm 0$，设“阅读”的中心端表示为$(1,0)$，“小王”与“通知”的依存端表示为$(4,0)$与$(3,0)$，则$S^{\mathrm{arc}}_{2,1}=4$、$S^{\mathrm{arc}}_{2,3}=3$，正好是Eisner例子中的两条弧分数。选定$2\to1$后，若关系头对nsubj、obj、其他的分数为$2,0,-1$，选nsubj。若两个中心词都正确却把$2\to1$标成obj，计入根弧的三条弧中UAS为1，LAS为$2/3$。

与Chen与Manning相比，双仿射解析把特征从固定窗口的嵌入换成整句的BiLSTM表示，把贪心的动作序列换成全局精确的树解码，代价是$O(n^2)$的弧分数和$O(n^3)$（或$O(n^2)$）的解码，不再是线性时间。Dozat与Manning在宾州树库的斯坦福依存版本上得到约95.7的UAS；其消融中去掉低维投影、直接用BiLSTM输出配对时LAS下降，而UAS几乎不变，支持了“关系判断需要专用表示”的解释\scite{nlp-analysis-dozat2017}。

\subsubsection{增强依存图预测}
\label{sec:nlp-analysis-enhanced}

基本树要求每个词只有一个中心词，但并列共享的论元和省略的谓词需要更多连接。\term{增强依存}（enhanced dependencies）在基本树之上按规范增加关系：允许一个词有多个中心词，允许在省略处插入空节点，某些结构中还可能出现环。它不是把树约束放宽后任意加边，哪些边成立仍由标注规范规定\scite{v9a-ud-enhanced}。

例如教学句“小王购买并阅读通知”，基本树有“购买”到“小王”的nsubj、“购买”到“阅读”的conj、“阅读”到“通知”的obj。增强表示把共享主语传播到第二个谓词，增加“阅读”到“小王”的nsubj；规范允许时还传播共享宾语，增加“购买”到“通知”的obj。“小王”因此有两条入边，这在增强图中是正确的。

增强图不能再用单头树解码，常见做法是把它当作多标签的弧预测。保留双仿射的评分结构，对允许的词对$(i,j)$和关系$r$计算$p_{i,j,r}=\operatorname{sigmoid}(S_{i,j,r})$，以增强弧的多热标签训练二元交叉熵，推断时保留超过阈值的弧。设新增两条弧的概率为$.85$与$.80$，反方向的错误候选为$.15$，验证集确定阈值$.5$，则保留前两条。阈值化不保证连通，可以再强制每个词至少保留一条最高分入边，或先用树解码得到基本树、再在其上加边。另一种做法完全不学习：读取基本树中的并列、控制与省略结构，按规范规则传播关系。

省略结构还需要空节点。对“甲买书，乙报纸”，第二分句缺少谓词“买”，增强表示在该处插入一个空节点，把“乙”“报纸”接到它下面。预测时要先确定空节点的位置与身份，再预测与它相连的弧；空节点不在原文中，不能把一个字符跨度强加给它。

评价改用增强弧指标。IWPT 2020共享任务采用ELAS，即增强图中“中心词、关系”对的F1，因为每个词的入边数不再固定为1，UAS与LAS的分母不再适用\scite{nlp-analysis-bouma2020}。表~\ref{tab:nlp-analysis-dependency-methods}按同一输出比较依存方法。转移系统学习动作，Eisner与Chu--Liu/Edmonds在给定弧分数上选树，双仿射网络负责产生弧分数；后三者常组合使用，并不互斥。

\begin{table*}[htbp]
  \centering
  \begin{tabularx}{\linewidth}{@{}l XXXX@{}}
    \toprule
    方法&解决的问题&核心机制&结构空间与复杂度&新代价\\
    \midrule
    arc-standard / arc-eager&线性时间解析&栈、缓冲区与连弧动作&投射树；$O(n)$，贪心或束&错误不可撤销；需交换动作处理非投射\\
    Eisner&投射树上的精确解码&四类区间状态的动态规划&投射树；$O(n^3)$&只支持一阶（或扩展状态后的高阶）分解\\
    Chu--Liu/Edmonds&非投射树的精确解码&贪心入边、收缩环、递归展开&任意生成树形图；$O(n^2)$至$O(n^3)$&失去投射约束的过滤；单根需额外处理\\
    双仿射&弧分数的学习&专用MLP投影与双仿射评分&接上述解码器；分数$O(n^2)$&需整句编码；非线性时间\\
    增强依存&共享论元与省略&多标签弧预测、空节点&一般有向图&无单头约束可用；需新评价口径\\
    \bottomrule
  \end{tabularx}
  \caption{依存解析方法对照。前三行是结构空间与搜索方式，第四行是评分函数，第五行改变了输出结构本身。}
  \label{tab:nlp-analysis-dependency-methods}
\end{table*}
